\documentclass[aspectratio=169]{beamer}
\usetheme{Madrid}
\usecolortheme{default}
\usepackage[utf8]{inputenc}
\usepackage[T5]{fontenc}
\usepackage{graphicx}
\usepackage{hyperref}
\usepackage{booktabs}

\setbeamertemplate{caption}[numbered]
\renewcommand{\figurename}{Hình}

\title[Trí tuệ Nhân tạo cho Kế toán]{Trí tuệ Nhân tạo cho Kế toán \\ \vspace{0.3cm} \Large KẾ HOẠCH SLIDE THỰC HÀNH - DAY 06 (BUỔI 6)}
\author{Đại học Đông Á}
\date{\today}

\begin{document}

% SLIDE 1
\begin{frame}
    \titlepage
\end{frame}

% SLIDE 2
\begin{frame}{Nội dung Bài học}
    \tableofcontents
\end{frame}


\section{KẾ HOẠCH SLIDE THỰC HÀNH - DAY 06 (BUỔI 6)}


\begin{frame}{Title: Buổi 6 TH: Thực hành Phân tích Báo cáo Tài chính.}
    \begin{itemize}
        \item Giảng viên: Đại học Đông Á
        \item Môn học: Trí tuệ Nhân tạo cho Kế toán
    \end{itemize}
\end{frame}


\begin{frame}{Năng lực đạt được sau buổi học.}
    \begin{itemize}
        \item \textbf{Về Kiến thức:} Nắm bắt quy trình phân tích dữ liệu đa chiều (Data Cubes) và tỷ số tài chính.
        \item \textbf{Về Kỹ năng:} Sử dụng ChatGPT (Advanced Data Analysis) để tính toán tự động và viết báo cáo đánh giá sức khỏe tài chính.
        \item \textbf{Về Tư duy:} Hình thành tư duy đối chiếu, kiểm chứng kết quả do AI tạo ra (Double-check), không phụ thuộc mù quáng.
    \end{itemize}
\end{frame}


\begin{frame}{Cấu trúc buổi thực hành (Agenda).}
    \begin{itemize}
        \item Chuẩn bị dữ liệu (Data Preparation).
        \item Khởi tạo môi trường AI \& Upload dữ liệu.
        \item Thực thi Prompt tính toán Tỷ số tài chính.
        \item Phân tích đa chiều (Pivot/Cubes) và trực quan hóa.
        \item Báo cáo đánh giá và phản biện AI.
        ### Phần 1: Chuẩn bị Dữ liệu (Data Preparation)
    \end{itemize}
\end{frame}


\begin{frame}{Giới thiệu bộ dữ liệu thực hành.}
    \begin{itemize}
        \item File Excel chứa Bảng Cân đối kế toán (Balance Sheet) và Báo cáo Kết quả Kinh doanh (Income Statement).
        \item Dữ liệu thô của 1 công ty niêm yết trong 3 năm liên tiếp.
    \end{itemize}
\end{frame}


\begin{frame}{Kiểm tra cấu trúc file Excel trước khi đưa vào AI.}
    \begin{itemize}
        \item Tại sao phải kiểm tra? Tránh lỗi "Garbage in, Garbage out".
        \item Các cột cần có: Khoản mục, Mã số, Năm T, Năm T-1, Năm T-2.
    \end{itemize}
\end{frame}


\begin{frame}{Làm sạch dữ liệu (Data Cleansing) cơ bản.}
    \begin{itemize}
        \item Xóa các dòng trống (Blank rows), cột thừa.
        \item Bỏ các định dạng phức tạp (Merge cells) có thể làm AI bối rối.
    \end{itemize}
\end{frame}


\begin{frame}{Chuẩn hóa tiêu đề cột (Headers).}
    \begin{itemize}
        \item Tiêu đề cột phải rõ ràng, ngắn gọn, nên dùng tiếng Anh hoặc tiếng Việt không dấu nếu công cụ AI dễ bị lỗi font.
        \item Dòng đầu tiên (Row 1) phải là tên cột.
    \end{itemize}
\end{frame}


\begin{frame}{Tầm quan trọng của định dạng số liệu (Number format).}
    \begin{itemize}
        \item Số tiền (VNĐ/USD) không được chứa ký tự chữ cái lẫn lộn.
        \item Phân biệt dấu phẩy (,) và dấu chấm (.) tùy theo Locale của Excel.
        ### Phần 2: Khởi tạo môi trường AI \& Upload Dữ liệu
    \end{itemize}
\end{frame}


\begin{frame}{Kích hoạt Advanced Data Analysis (ADA) trên ChatGPT.}
    \begin{itemize}
        \item Đăng nhập tài khoản ChatGPT.
        \item Đảm bảo tính năng phân tích dữ liệu (Data Analyst) được bật.
    \end{itemize}
\end{frame}


\begin{frame}{Thao tác Upload file an toàn.}
    \begin{itemize}
        \item Kéo thả file Excel vào khung chat.
        \item Lưu ý bảo mật: Đổi tên công ty thực thành công ty giả định (Công ty ABC) để bảo vệ dữ liệu nội bộ.
    \end{itemize}
\end{frame}


\begin{frame}{Prompt số 1 - Lệnh "Đọc hiểu" dữ liệu.}
    \begin{itemize}
        \item "Đóng vai chuyên gia phân tích tài chính. Hãy đọc file Excel đính kèm và tóm tắt cấu trúc dữ liệu của 2 sheet: Balance Sheet và Income Statement."
    \end{itemize}
\end{frame}


\begin{frame}{Phân tích phản hồi đầu tiên của AI.}
    \begin{itemize}
        \item Kiểm tra xem AI có nhận diện đúng số dòng, số cột và các khoản mục chính không.
        \item Nếu AI đọc sai (nhận nhầm header), cần nhắc nhở và sửa lỗi ngay.
        ### Phần 3: Thực thi Prompt Tính toán Tỷ số tài chính
    \end{itemize}
\end{frame}


\begin{frame}{Prompt số 2 - Tính toán Nhóm Thanh khoản (Liquidity).}
    \begin{itemize}
        \item "Dựa trên dữ liệu, hãy tạo một bảng tính Tỷ số thanh toán hiện hành và Tỷ số thanh toán nhanh cho cả 3 năm."
    \end{itemize}
\end{frame}


\begin{frame}{Đọc kết quả Tỷ số Thanh khoản.}
    \begin{itemize}
        \item So sánh kết quả AI trả về với công thức chuẩn.
        \item Kiểm tra xem AI lấy số liệu "Tài sản ngắn hạn" ở dòng nào.
    \end{itemize}
\end{frame}


\begin{frame}{Prompt số 3 - Tính toán Nhóm Đòn bẩy (Solvency).}
    \begin{itemize}
        \item "Tiếp tục tính Tỷ số Nợ / Vốn chủ sở hữu (D/E) và Tỷ số Nợ / Tổng tài sản cho 3 năm."
    \end{itemize}
\end{frame}


\begin{frame}{Đọc kết quả Tỷ số Đòn bẩy.}
    \begin{itemize}
        \item Phân tích rủi ro từ con số AI đưa ra.
        \item Xác minh AI có cộng đúng tổng nợ (Nợ ngắn + Nợ dài) hay không.
    \end{itemize}
\end{frame}


\begin{frame}{Prompt số 4 - Tính toán Nhóm Sinh lời (Profitability).}
    \begin{itemize}
        \item "Hãy tính ROA, ROE và Biên lợi nhuận gộp. Trình bày dưới dạng bảng so sánh."
    \end{itemize}
\end{frame}


\begin{frame}{Đọc kết quả Tỷ số Sinh lời.}
    \begin{itemize}
        \item AI thường tính ROA bằng Tổng tài sản cuối kỳ hay Bình quân? (Yêu cầu AI giải thích công thức nó đã dùng).
    \end{itemize}
\end{frame}


\begin{frame}{Prompt nâng cao - Mô hình DuPont.}
    \begin{itemize}
        \item "Sử dụng dữ liệu hiện có, hãy phân tách ROE của năm gần nhất theo Mô hình DuPont 3 nhân tố."
    \end{itemize}
\end{frame}


\begin{frame}{Phân tích kết quả DuPont do AI xuất ra.}
    \begin{itemize}
        \item Nhìn vào 3 nhân tố (Biên lợi nhuận ròng, Vòng quay TS, Đòn bẩy).
        \item Nguyên nhân nào làm thay đổi ROE?
        ### Phần 4: Phân tích Đa chiều (Từ Spreadsheet sang Data Cubes)
    \end{itemize}
\end{frame}


\begin{frame}{Giới thiệu khái niệm Data Cubes cơ bản.}
    \begin{itemize}
        \item Vượt ra khỏi bảng 2 chiều (Dòng x Cột).
        \item Thêm các chiều phân tích (Dimension): Thời gian, Chi nhánh, Dòng sản phẩm.
    \end{itemize}
\end{frame}


\begin{frame}{Data Cubes trong AI.}
    \begin{itemize}
        \item Khác với Excel Pivot Table phải kéo thả thủ công.
        \item AI có thể tạo các lát cắt dữ liệu (Slice and Dice) bằng mã Python ngầm định.
    \end{itemize}
\end{frame}


\begin{frame}{Prompt số 5 - Tạo Data Cubes giả lập.}
    \begin{itemize}
        \item "Hãy nhóm doanh thu và chi phí theo từng quý (nếu có dữ liệu) và so sánh sự tăng trưởng (YoY)."
    \end{itemize}
\end{frame}


\begin{frame}{Trực quan hóa dữ liệu (Data Visualization).}
    \begin{itemize}
        \item Yêu cầu AI: "Hãy vẽ biểu đồ cột thể hiện sự biến động của Doanh thu thuần và Lợi nhuận sau thuế trong 3 năm."
    \end{itemize}
\end{frame}


\begin{frame}{Trực quan hóa cấu trúc vốn.}
    \begin{itemize}
        \item Yêu cầu AI: "Vẽ biểu đồ tròn (Pie chart) thể hiện cơ cấu Nợ và Vốn chủ sở hữu của năm mới nhất."
    \end{itemize}
\end{frame}


\begin{frame}{Trực quan hóa xu hướng tỷ số.}
    \begin{itemize}
        \item Yêu cầu AI: "Vẽ biểu đồ đường (Line chart) mô tả xu hướng của ROA và ROE, thêm đường trung bình ngành (giả định là 10\%)."
        ### Phần 5: Báo cáo Đánh giá và Phản biện AI
    \end{itemize}
\end{frame}


\begin{frame}{Prompt Tổng hợp - Viết báo cáo Sức khỏe tài chính.}
    \begin{itemize}
        \item "Dựa trên tất cả các tỷ số đã tính, hãy viết một đoạn nhận xét 300 chữ đánh giá sức khỏe tài chính của công ty. Nêu rõ 2 điểm mạnh và 2 rủi ro."
    \end{itemize}
\end{frame}


\begin{frame}{Đánh giá Báo cáo của AI.}
    \begin{itemize}
        \item Báo cáo có logic không? Các nhận định có khớp với số liệu không?
        \item Phát hiện những lời "sáo rỗng" hoặc "ảo giác" của AI.
    \end{itemize}
\end{frame}


\begin{frame}{Phản biện lại AI.}
    \begin{itemize}
        \item Kỹ năng quan trọng: Yêu cầu AI giải thích "Tại sao bạn kết luận rủi ro thanh khoản cao khi Tỷ số hiện hành là 1.5?".
        \item Xem AI giải thích các giả định của nó.
    \end{itemize}
\end{frame}


\begin{frame}{Kết xuất kết quả cuối cùng.}
    \begin{itemize}
        \item Hướng dẫn sinh viên copy/paste báo cáo và biểu đồ từ ChatGPT vào Word/PowerPoint để hoàn thiện bài tập nhóm.
    \end{itemize}
\end{frame}


\begin{frame}{Tổng kết buổi thực hành.}
    \begin{itemize}
        \item Điểm lại các lệnh Prompt hiệu quả nhất đã sử dụng.
        \item Nhắc nhở: AI chỉ là công cụ hỗ trợ tính toán nhanh, linh hồn của bài phân tích vẫn là tư duy tài chính của người làm kế toán.
    \end{itemize}
\end{frame}


\begin{frame}{Q\&A - Giải đáp thắc mắc.}
    \begin{itemize}
        \item Sinh viên đặt câu hỏi về các lỗi gặp phải khi thực thi Prompt.
        \item Bài tập về nhà: Tự tải BCTC của 1 công ty khác (VNM, FPT) và thực hiện quy trình tương tự.
    \end{itemize}
\end{frame}


\end{document}